\documentclass[11pt,a4paper]{article}

% Paquetes esenciales
\usepackage[english]{babel}
\usepackage[margin=1in]{geometry}
\usepackage{graphicx}
\usepackage{amsmath}
\usepackage{amssymb}
\usepackage{hyperref}
\usepackage{cite}
\usepackage{natbib}
\usepackage{booktabs}
\usepackage{longtable}
\usepackage{xcolor}
\usepackage{listings}
\usepackage{fancyhdr}
\usepackage{lastpage}
\usepackage{float}
\usepackage{array}
\usepackage{multirow}

% Configuración de listings para código
\lstset{
    language=Python,
    basicstyle=\ttfamily\small,
    breaklines=true,
    breakatwhitespace=true,
    keywordstyle=\color{blue},
    commentstyle=\color{gray},
    stringstyle=\color{red},
    backgroundcolor=\color{gray!10},
    frame=single,
    captionpos=b,
    numbers=left,
    numberstyle=\tiny
}
\setlength{\parskip}{0.6em}

% Header y Footer
\pagestyle{fancy}
\fancyhf{}
\lhead{Recipe Recommendation System}
\chead{Week 3 Milestone}
\rhead{Semester Group Project}
\cfoot{\thepage\ of \pageref{LastPage}}

\begin{document}

\begin{titlepage}
    \begin{center}


        % Si luego quieres agregar logo, va aquí:
         \includegraphics[width=3.5cm]{logo-upc-png-transparente-1.png}
         \vspace{1cm}

        {\Large \textbf{UNIVERSIDAD PERUANA DE CIENCIAS APLICADAS}}\\[1.2cm]

        {\Large \textbf{FALCULTY OF ENGINEERING}}\\[1.8cm]

        {\Large \textbf{PROJECT REPORT - RECIPE RECOMMENDATIONS PROJECT}}\\[1.2cm]

        {\large Subject: Big Data}\\[0.3cm]
        {\large NRC: 18517}\\[1.5cm]

        {\large \textbf{Team members:}}\\[0.8cm]

        \renewcommand{\arraystretch}{1.6}
        \begin{tabular}{|>{\centering\arraybackslash}m{7cm}|>{\centering\arraybackslash}m{3cm}|>{\centering\arraybackslash}m{3cm}|}
            \hline
            \textbf{Full Name} & \textbf{Student ID} \\
            \hline
            Marchan Quispe, Mildred Micaela&U202213292 \\
            \hline
            Medina Agnini, Bruno Alessandro&U202216661 \\
            \hline
            Miranda Rafael, Nicolás Francisco&U202216895 \\
            \hline
            Rodríguez Valencia, Rosa Maria&U202212675 \\
            \hline
        \end{tabular}

        \vfill

        {\large \textbf{Professor:} Alarcon Delgado Carlos Adrian}\\[1.5cm]

        {\LARGE 2026 - 01}

    \end{center}
\end{titlepage}



\newpage
\tableofcontents
\newpage
% ============================================================================
% SECTION 1: PROJECT PROPOSAL
% ============================================================================
\section{Project Proposal}

\subsection{Domain and Problem Context}

Food discovery and recipe recommendation is a high-impact domain with significant real-world applications. Users regularly struggle with meal planning, ingredient utilization, and finding recipes aligned with their preferences and dietary constraints. Current platforms offer limited guidance on how recipes relate to one another through shared ingredients, cooking techniques, or flavor profiles.

\textbf{Domain Focus:} Online recipe curation, discovery, and personalized recommendation.

\textbf{Core Problem:} Given a corpus of recipes with rich metadata (ingredients, instructions, ratings, cook time, difficulty), build a system that:
\begin{enumerate}
    \item Segments recipes into meaningful clusters based on ingredient, flavor, and technique similarity.
    \item Recommends recipes to users based on their interaction history (views, ratings, saves).
    \item Uncovers structural patterns through an ingredient-ingredient co-occurrence network.
    \item Ranks ingredients and recipes by centrality and influence.
\end{enumerate}

\subsection{Expected Product Questions}

The system will answer questions such as:

\begin{itemize}
    \item \textbf{Discovery:} Given a recipe I enjoyed, what other recipes might I like?
    \item \textbf{Segmentation:} Are there natural clusters of recipes (e.g., quick weeknight dinners, elaborate desserts, regional cuisines)?
    \item \textbf{Ingredient Insights:} Which ingredients are central to the recipe ecosystem, and which recipes share them?
    \item \textbf{Ranking:} Which recipes are most highly rated? Which ingredients appear most frequently in top-rated recipes?
    \item \textbf{Connectivity:} How does the ingredient network reveal hidden relationships between recipes?
\end{itemize}

\subsection{Why This Dataset Supports the Full Course}

This project satisfies the second-half requirements because:

\begin{enumerate}
    \item \textbf{Feature Richness:} Recipes offer numeric (ratings, cook time, difficulty), categorical (cuisine, diet type), and text features (instructions, ingredient descriptions).
    \item \textbf{Interaction Layer:} User-recipe interactions (views, ratings, saves, reviews) create a rich behavioral signal.
    \item \textbf{Graph Potential:} The ingredient network naturally induces a bipartite graph (recipes $\leftrightarrow$ ingredients), which can be projected into item-item or ingredient-ingredient spaces.
    \item \textbf{Dimensionality Challenge:} The ingredient matrix (recipes $\times$ ingredients) is sparse and high-dimensional, necessitating PCA, SVD, or embedding methods.
    \item \textbf{Clustering and Ranking:} Both unsupervised segmentation and supervised or collaborative ranking are meaningful on this data.
    \item \textbf{Non-Trivial Scale:} Public recipe datasets contain 10k--100k+ recipes with thousands of unique ingredients, sufficient for meaningful analysis.
\end{enumerate}

\newpage

% ============================================================================
% SECTION 2: SOURCE INVENTORY
% ============================================================================
\section{Source Inventory}

\subsection{Primary Data Sources}

\subsubsection{Kaggle Food.com Recipes and Reviews Dataset}

\begin{table}[H]
    \centering
    \begin{tabular}{|p{3cm}|p{10cm}|}
        \hline
        \textbf{Attribute} & \textbf{Value} \\
        \hline
        \textbf{Source} & Kaggle -- ``Food.com Recipes and Reviews'' \\
        \hline
        \textbf{URL} & \url{https://www.kaggle.com/datasets/irkaal/foodcom-recipes-and-reviews} \\
        \hline
        \textbf{License} & CC0: Public Domain (verified on dataset page) \\
        \hline
        \textbf{Access} & Public download, no authentication required \\
        \hline
        \textbf{Raw Format} & CSV files (recipes.csv, interactions.csv, and related tables) \\
        \hline
        \textbf{Estimated Size} & 100--300 MB uncompressed \\
        \hline
        \textbf{Records} & ~200,000 recipes; ~1,000,000+ reviews/interactions; ~10,000+ unique ingredients \\
        \hline
    \end{tabular}
    \caption{Primary Data Source: Kaggle Food.com Recipes and Reviews Dataset}
    \label{tab:primary_source}
\end{table}

\subsection{Alternative or Supplementary Sources}

Should the primary Kaggle source become unavailable, the team will pivot to:

\begin{itemize}
    \item \textbf{Food.com Recipe Dumps:} Publicly available recipe and review data from the Food.com platform via research releases.
    \item \textbf{AllRecipes Scrape (Historical):} Ethically obtained historical recipe corpora from academic repositories.
    \item \textbf{Recipe1M Dataset:} Large-scale multilingual recipe dataset from a published academic benchmark.
\end{itemize}

\subsection{License and Legal Terms}

\begin{table}[H]
    \centering
    \begin{tabular}{|p{4cm}|p{9cm}|}
        \hline
        \textbf{Attribute} & \textbf{Details} \\
        \hline
        \textbf{License Type} & CC0 1.0 Universal (Public Domain Dedication) \\
        \hline
        \textbf{Permitted Uses} & Unrestricted use, modification, distribution for academic and commercial purposes \\
        \hline
        \textbf{Attribution} & Optional (CC0 does not require attribution, but team will credit the source) \\
        \hline
        \textbf{Restrictions} & None; data is in the public domain \\
        \hline
    \end{tabular}
    \caption{License and Access Conditions}
    \label{tab:license}
\end{table}

\newpage

% ============================================================================
% SECTION 3: SCHEMA DRAFT
% ============================================================================
\section{Schema Draft}

\subsection{Entity-Relationship Diagram Description}

The system models recipes, ingredients, and user interactions. The schema is normalized to support joins and graph construction.

\subsection{Core Entity Tables}

\subsubsection{recipes}

Primary entity table.

\begin{table}[H]
    \centering
    \begin{tabular}{|p{2.5cm}|p{2cm}|p{8cm}|}
        \hline
        \textbf{Column} & \textbf{Type} & \textbf{Description} \\
        \hline
        recipe\_id & INTEGER & Unique recipe identifier (Primary Key) \\
        \hline
        title & TEXT & Recipe name \\
        \hline
        cuisine & VARCHAR(50) & Cuisine category (e.g., Italian, Asian, American) \\
        \hline
        diet\_type & VARCHAR(50) & Dietary category (e.g., vegetarian, vegan, gluten-free) \\
        \hline
        cook\_time & INTEGER & Cooking time in minutes \\
        \hline
        prep\_time & INTEGER & Preparation time in minutes \\
        \hline
        difficulty & VARCHAR(20) & Difficulty level (easy, medium, hard) \\
        \hline
        instructions & TEXT & Step-by-step cooking instructions \\
        \hline
        avg\_rating & FLOAT & Mean user rating (0--5 scale) \\
        \hline
        num\_ratings & INTEGER & Count of user ratings \\
        \hline
        created\_date & DATE & Date recipe was added \\
        \hline
    \end{tabular}
    \caption{recipes Table Schema}
    \label{tab:recipes_schema}
\end{table}

\subsubsection{ingredients}

Ingredient reference table.

\begin{table}[H]
    \centering
    \begin{tabular}{|p{2.5cm}|p{2cm}|p{8cm}|}
        \hline
        \textbf{Column} & \textbf{Type} & \textbf{Description} \\
        \hline
        ingredient\_id & INTEGER & Unique ingredient identifier (Primary Key) \\
        \hline
        name & VARCHAR(100) & Standardized ingredient name \\
        \hline
        category & VARCHAR(50) & Ingredient category (e.g., vegetable, spice, grain) \\
        \hline
        unit & VARCHAR(20) & Standard unit of measure (g, ml, cup, tbsp) \\
        \hline
    \end{tabular}
    \caption{ingredients Table Schema}
    \label{tab:ingredients_schema}
\end{table}

\subsubsection{recipe\_ingredients}

Junction table linking recipes to ingredients (Many-to-Many relationship).

\begin{table}[H]
    \centering
    \begin{tabular}{|p{2.5cm}|p{2cm}|p{8cm}|}
        \hline
        \textbf{Column} & \textbf{Type} & \textbf{Description} \\
        \hline
        recipe\_id & INTEGER & Foreign Key to recipes (Primary Key Part 1) \\
        \hline
        ingredient\_id & INTEGER & Foreign Key to ingredients (Primary Key Part 2) \\
        \hline
        quantity & FLOAT & Numeric quantity \\
        \hline
        unit & VARCHAR(20) & Measurement unit \\
        \hline
    \end{tabular}
    \caption{recipe\_ingredients Table Schema}
    \label{tab:recipe_ingredients_schema}
\end{table}

\subsubsection{ratings}

User interaction table (ratings and reviews).

\begin{table}[H]
    \centering
    \begin{tabular}{|p{2.5cm}|p{2cm}|p{8cm}|}
        \hline
        \textbf{Column} & \textbf{Type} & \textbf{Description} \\
        \hline
        rating\_id & INTEGER & Unique rating identifier (Primary Key) \\
        \hline
        recipe\_id & INTEGER & Foreign Key to recipes \\
        \hline
        user\_id & INTEGER & Unique user identifier \\
        \hline
        score & FLOAT & Numeric rating (0--5 scale) \\
        \hline
        review\_text & TEXT & User review comment (optional) \\
        \hline
        date & DATE & Date of rating \\
        \hline
    \end{tabular}
    \caption{ratings Table Schema}
    \label{tab:ratings_schema}
\end{table}

\subsection{Expected Joins}

\begin{itemize}
    \item \textbf{recipes \& ratings:} One-to-Many join on recipe\_id (one recipe has many ratings).
    \item \textbf{recipes \& ingredients:} Many-to-Many join via recipe\_ingredients junction table.
    \item \textbf{ratings \& recipes:} Used to aggregate user behavior by recipe and to build item-item similarity.
\end{itemize}

\subsection{Graph Construction Sketch}

\textbf{Ingredient-Ingredient Network:} 
\begin{equation}
\text{Edge}(i_a, i_b) = \text{number of recipes containing both } i_a \text{ and } i_b
\end{equation}

This creates an undirected, weighted graph of ingredient co-occurrences, enabling discovery of complementary ingredients and regional flavor patterns.

\textbf{Recipe-Recipe Network (later)}: Projected from ingredients or built from rating co-occurrence.

\newpage

% ============================================================================
% SECTION 4: PROCESSED DATASET V1
% ============================================================================
\section{Processed Dataset V1}

\subsection{Data Ingestion and Cleaning Pipeline}

The team will implement a reproducible Python ingestion script (using pandas, SQLite, or polars) that:

\begin{enumerate}
    \item Downloads or reads the raw CSV files from the source.
    \item Removes or flags duplicate recipes (by title + ingredient set hash).
    \item Standardizes ingredient names (lowercasing, removing punctuation, merging synonyms).
    \item Handles missing values in ratings, cook time, and dietary categories.
    \item Validates foreign key relationships.
    \item Saves cleaned tables to the processed/ directory.
\end{enumerate}

\subsection{Data Cleaning Decisions}

\begin{table}[H]
    \centering
    \begin{tabular}{|p{2.5cm}|p{8cm}|}
        \hline
        \textbf{Issue} & \textbf{Handling Strategy} \\
        \hline
        Missing cook time & Impute with median by cuisine; flag as ``unknown'' \\
        \hline
        Duplicate recipes & Remove exact duplicates; merge near-duplicates by ingredient set \\
        \hline
        Non-standard ingredient names & Normalize: lowercase, remove extra spaces, merge variants \\
        \hline
        Missing ratings & Retain recipes; use recipe-level aggregates (avg\_rating) \\
        \hline
        Outliers (e.g., 300-hour cook time) & Flag and inspect; cap at 95th percentile if errors \\
        \hline
        NULL dietary categories & Infer from ingredient keywords or label as ``unspecified'' \\
        \hline
    \end{tabular}
    \caption{Data Cleaning Decisions}
    \label{tab:cleaning}
\end{table}

\subsection{Output Artifacts}

After processing, the team will produce:

\begin{lstlisting}[language=bash, caption=Output Files from Ingestion Pipeline]
data/processed/
  recipes_clean.csv          # Main recipe table
  ingredients_clean.csv      # Ingredient reference
  recipe_ingredients.csv     # Junction table
  ratings_clean.csv          # User interaction data
  data_dictionary.md         # Full metadata
\end{lstlisting}

\subsection{Data Quality Checks}

The ingestion script will output a quality report:

\begin{table}[H]
    \centering
    \begin{tabular}{|p{3cm}|p{3cm}|p{4.5cm}|}
        \hline
        \textbf{Metric} & \textbf{Method} & \textbf{Interpretation} \\
        \hline
        Duplicate recipes & Hash(title + ingredients) & Count and remove \\
        \hline
        Foreign key violations & Check recipe\_id in recipe\_ingredients & Count unmatched \\
        \hline
        Missing ingredient names & NULL check in ingredients.name & Report count \\
        \hline
        Invalid ratings & $\text{score} \notin [0, 5]$ & Flag and remove \\
        \hline
    \end{tabular}
    \caption{Data Quality Checks}
    \label{tab:quality_checks}
\end{table}

\newpage

% ============================================================================
% SECTION 5: DATA DICTIONARY DRAFT
% ============================================================================
\section{Data Dictionary Draft}

\subsection{recipes\_clean.csv}

\begin{longtable}{|p{2.5cm}|p{1.5cm}|p{2.5cm}|p{6.5cm}|}
    \hline
    \textbf{Column} & \textbf{Type} & \textbf{Non-Null \%} & \textbf{Description and Range} \\
    \hline
    \endhead
    
    recipe\_id & int64 & 100 & Unique identifier. Range: [1, n\_recipes]. \\
    \hline
    title & object & 100 & Recipe name (string). Max length: 255 characters. \\
    \hline
    cuisine & object & 85 & Categorical cuisine label (string). Common values: Italian, Asian, American, Mediterranean, Mexican. \\
    \hline
    diet\_type & object & 70 & Dietary classification (string). Values: vegetarian, vegan, gluten-free, dairy-free, paleo, unspecified. \\
    \hline
    cook\_time & float64 & 90 & Time in minutes. Range: [0, 300]. Mean: ~30 min. \\
    \hline
    prep\_time & float64 & 88 & Preparation time in minutes. Range: [0, 120]. Mean: ~15 min. \\
    \hline
    difficulty & object & 92 & Categorical difficulty. Values: easy, medium, hard. \\
    \hline
    instructions & object & 99 & Full text of cooking steps (string). Avg length: 200--500 words. \\
    \hline
    avg\_rating & float64 & 95 & Mean user rating. Range: [1.0, 5.0]. Median: ~4.2. \\
    \hline
    num\_ratings & int64 & 100 & Count of ratings per recipe. Range: [1, 500]. Median: ~20. \\
    \hline
    created\_date & object & 80 & ISO date (YYYY-MM-DD format). Span: 2010--2023. \\
    \hline
    
    \caption{Data Dictionary: recipes\_clean.csv}
    \label{tab:dd_recipes}
\end{longtable}

\subsection{ingredients\_clean.csv}

\begin{longtable}{|p{2.5cm}|p{1.5cm}|p{2.5cm}|p{6.5cm}|}
    \hline
    \textbf{Column} & \textbf{Type} & \textbf{Non-Null \%} & \textbf{Description and Range} \\
    \hline
    \endhead
    
    ingredient\_id & int64 & 100 & Unique identifier. Range: [1, n\_ingredients]. \\
    \hline
    name & object & 100 & Standardized ingredient name. Max length: 100 characters. All lowercase. \\
    \hline
    category & object & 88 & Ingredient category (string). Values: vegetable, spice, grain, protein, fruit, dairy, oil, condiment. \\
    \hline
    unit & object & 70 & Standard unit. Values: g, ml, cup, tbsp, tsp, oz, lb, count. \\
    \hline
    
    \caption{Data Dictionary: ingredients\_clean.csv}
    \label{tab:dd_ingredients}
\end{longtable}

\subsection{recipe\_ingredients.csv}

\begin{longtable}{|p{2.5cm}|p{1.5cm}|p{2.5cm}|p{6.5cm}|}
    \hline
    \textbf{Column} & \textbf{Type} & \textbf{Non-Null \%} & \textbf{Description and Range} \\
    \hline
    \endhead
    
    recipe\_id & int64 & 100 & Foreign key to recipes. \\
    \hline
    ingredient\_id & int64 & 100 & Foreign key to ingredients. \\
    \hline
    quantity & float64 & 95 & Numeric quantity. Range: [0.01, 100]. \\
    \hline
    unit & object & 95 & Measurement unit for the quantity. \\
    \hline
    
    \caption{Data Dictionary: recipe\_ingredients.csv}
    \label{tab:dd_recipe_ingredients}
\end{longtable}

\subsection{ratings\_clean.csv}

\begin{longtable}{|p{2.5cm}|p{1.5cm}|p{2.5cm}|p{6.5cm}|}
    \hline
    \textbf{Column} & \textbf{Type} & \textbf{Non-Null \%} & \textbf{Description and Range} \\
    \hline
    \endhead
    
    rating\_id & int64 & 100 & Unique identifier for each rating. \\
    \hline
    recipe\_id & int64 & 100 & Foreign key to recipes. \\
    \hline
    user\_id & int64 & 100 & Unique user identifier (anonymized). \\
    \hline
    score & float64 & 100 & Numeric rating. Range: [0, 5]. \\
    \hline
    review\_text & object & 40 & Optional user review comment. Max length: 1000 characters. \\
    \hline
    date & object & 95 & ISO date (YYYY-MM-DD format). Span: 2010--2023. \\
    \hline
    
    \caption{Data Dictionary: ratings\_clean.csv}
    \label{tab:dd_ratings}
\end{longtable}

\newpage

% ============================================================================
% SECTION 6: SCALE ANALYSIS
% ============================================================================
\section{Scale Analysis}

\subsection{Data Volume Estimates}

\begin{table}[H]
    \centering
    \begin{tabular}{|p{3.5cm}|p{2.5cm}|p{2.5cm}|p{3cm}|}
        \hline
        \textbf{Table} & \textbf{Rows} & \textbf{Columns} & \textbf{Est. Size} \\
        \hline
        recipes & 200,000 & 12+ & ~50 MB \\
        \hline
        interactions/ratings & 1,000,000 & 5+ & ~100 MB \\
        \hline
        ingredients & 10,000 & 4 & ~2 MB \\
        \hline
        recipe\_ingredients & 500,000 & 4 & ~30 MB \\
        \hline
        \textbf{Total} & \textbf{1,710,000+} & & \textbf{~182 MB} \\
        \hline
    \end{tabular}
    \caption{Estimated Data Scale}
    \label{tab:scale}
\end{table}

\subsection{Sparsity and Dimensionality}

\subsubsection{Recipe-Ingredient Matrix}

If we construct a binary or count matrix $R \in \{0, 1, \ldots\}^{n_{\text{recipes}} \times n_{\text{ingredients}}}$, where $R_{ij} = 1$ if recipe $i$ contains ingredient $j$:

\begin{equation}
\text{Density} = \frac{\text{Non-zero entries}}{n_{\text{recipes}} \times n_{\text{ingredients}}} = \frac{500,000}{200,000 \times 10,000} \approx 0.025\%
\end{equation}

This is an extremely sparse matrix, even more sparse than typical collaborative filtering datasets, which is perfect for advanced dimensionality reduction techniques (PCA, SVD, autoencoders).

\subsubsection{User-Recipe Matrix}

If we construct a user-recipe interaction matrix $I \in \mathbb{R}^{n_{\text{users}} \times n_{\text{recipes}}}$ where $I_{ui}$ is the rating:

\begin{equation}
\text{Est. unique users} \approx \sqrt{\text{total interactions} \times 0.3} \approx 20,000 \text{ to } 50,000
\end{equation}

\begin{equation}
\text{Density} = \frac{1,000,000}{35,000 \times 200,000} \approx 0.14\%
\end{equation}

Still sparse, but with much larger scale. This provides excellent opportunity for collaborative filtering and matrix factorization techniques.

\subsection{Missing Data Profile}

\begin{table}[H]
    \centering
    \begin{tabular}{|p{3cm}|p{2cm}|p{2cm}|p{4cm}|}
        \hline
        \textbf{Column} & \textbf{Table} & \textbf{Missing \%} & \textbf{Action} \\
        \hline
        cuisine & recipes & 15 & Impute with ``unspecified'' or infer from ingredients \\
        \hline
        diet\_type & recipes & 30 & Infer from ingredients or leave as ``unspecified'' \\
        \hline
        cook\_time & recipes & 10 & Impute with median by cuisine \\
        \hline
        review\_text & ratings & 60 & Optional; keep NULL for non-text analysis \\
        \hline
        created\_date & recipes & 20 & Impute with median date by cuisine \\
        \hline
    \end{tabular}
    \caption{Missing Data Profile and Handling}
    \label{tab:missing}
\end{table}

\subsection{Data Growth Projection}

If the team adds recipes over time or extends to multiple recipe platforms, the scale may grow to:

\begin{equation}
\text{Recipes: } 200,000 \to 500,000+; \quad \text{Interactions: } 1,000,000 \to 5,000,000+
\end{equation}

The Food.com dataset is already at substantial scale. The team will design the pipeline to scale incrementally, using batch processing and efficient indexing.

\newpage

% ============================================================================
% SECTION 7: ETHICS AND ACCESS NOTE
% ============================================================================
\section{Ethics and Access Note}

\subsection{Data Provenance and Legality}

The Recipe Recommendation project uses \textbf{publicly available, anonymized data} from Kaggle, released under the \textbf{CC0 Public Domain} license. The data consists of:

\begin{itemize}
    \item Recipe metadata (title, ingredients, instructions, cook time, difficulty).
    \item Aggregated user ratings and reviews.
    \item No personally identifiable information (names, email addresses, IP addresses) is retained.
    \item User IDs are anonymized (integer hashes with no mapping to real users).
\end{itemize}

\subsection{Justification for Use}

\begin{enumerate}
    \item \textbf{Public Domain:} The dataset is released under CC0, permitting unrestricted academic and commercial use.
    \item \textbf{No Scraping:} Data is obtained via official Kaggle download, not unauthorized scraping of private platforms.
    \item \textbf{Anonymized Users:} User IDs are already anonymized in the source dataset; no re-identification is possible.
    \item \textbf{Permitted Aggregation:} Aggregating ratings and building item-item graphs does not violate individual privacy.
\end{enumerate}

\subsection{Personal Data Risks}

Potential risks and mitigations:

\begin{table}[H]
    \centering
    \begin{tabular}{|p{3cm}|p{4.5cm}|p{4.5cm}|}
        \hline
        \textbf{Risk} & \textbf{Potential Harm} & \textbf{Mitigation} \\
        \hline
        Re-identification & User IDs could theoretically be linked to real users via external data & Use integer hashes; never merge with external user databases \\
        \hline
        Inference attacks & Patterns in recipes + ratings could reveal user preferences & Aggregate at item level; no user-level outputs \\
        \hline
        Data leakage & Raw data could be misused if left unsecured & Store in locked project folder; no public data dumps \\
        \hline
        Bias in ratings & Low-income or non-English recipes may be underrepresented & Document dataset bias; note limitations in final report \\
        \hline
    \end{tabular}
    \caption{Personal Data Risks and Mitigations}
    \label{tab:risks}
\end{table}

\subsection{Data Handling Practices}

The team commits to:

\begin{enumerate}
    \item \textbf{Anonymity:} Never attempt to re-identify users or retain personal data beyond anonymized IDs.
    \item \textbf{Minimal Processing:} Only use data for this course project; no secondary use or redistribution.
    \item \textbf{Secure Storage:} Keep processed data on password-protected machines or institutional servers.
    \item \textbf{Transparency:} Document all data-handling decisions in the project README and final report.
    \item \textbf{Compliance:} Comply with Kaggle Terms of Service and any platform-specific restrictions.
\end{enumerate}

\subsection{Ethics Approval}

This project does not require Institutional Review Board (IRB) approval because:

\begin{itemize}
    \item Data is already publicly available and anonymized.
    \item No human subjects are directly contacted or surveyed.
    \item Analysis is purely observational on aggregated recipe and rating data.
    \item No sensitive attributes (health, race, gender, location) are used in analysis.
\end{itemize}

\newpage

% ============================================================================
% SECTION 8: TECHNICAL SPECIFICATION FOR REPRODUCIBILITY
% ============================================================================
\section{Technical Specification for Reproducibility}

\subsection{Repository Structure}

The team will organize the project as:

\begin{lstlisting}[language=bash]
recipe-recommendation-system/
├── README.md
├── requirements.txt
├── data/
│   ├── raw/
│   │   ├── recipes.csv
│   │   ├── ingredients.csv
│   │   ├── recipe_ingredients.csv
│   │   └── ratings.csv
│   ├── interim/
│   │   └── (intermediate processing files)
│   └── processed/
│       ├── recipes_clean.csv
│       ├── ingredients_clean.csv
│       ├── recipe_ingredients.csv
│       ├── ratings_clean.csv
│       └── data_dictionary.md
├── src/
│   ├── __init__.py
│   ├── ingest.py      # Ingestion & cleaning script
│   ├── features.py    # Feature engineering
│   ├── clustering.py  # Clustering methods
│   ├── recommend.py   # Recommendation engine
│   └── graph.py       # Graph construction
├── notebooks/
│   ├── 01_eda.ipynb
│   └── 02_scale_analysis.ipynb
├── artifacts/
│   └── (outputs: matrices, embeddings, graphs, etc.)
├── reports/
│   ├── week3_report.pdf
│   ├── week5_report.pdf
│   └── ...
└── scripts/
    ├── build.sh       # End-to-end pipeline
    └── validate.py    # Data quality checks
\end{lstlisting}

\subsection{Ingestion Command}

The team will document a single reproducible command:

\begin{lstlisting}[language=bash]
# Step 1: Download raw data (manual or automated)
python src/ingest.py --input data/raw/ --output data/processed/

# Step 2: Validate outputs
python scripts/validate.py --data data/processed/
\end{lstlisting}

\subsection{Dependencies}

Required Python packages:

\begin{lstlisting}[caption=requirements.txt]
pandas>=1.3.0
numpy>=1.21.0
scikit-learn>=1.0.0
scipy>=1.7.0
networkx>=2.6.0
matplotlib>=3.4.0
seaborn>=0.11.0
jupyter>=1.0.0
\end{lstlisting}

\subsection{Expected Outputs}

After running the ingestion pipeline:

\begin{enumerate}
    \item Four CSV files in data/processed/ (recipes, ingredients, recipe\_ingredients, ratings).
    \item A data\_dictionary.md file describing all columns.
    \item A validation report (JSON or log) listing row counts, missing data, and key statistics.
\end{enumerate}

\newpage

% ============================================================================
% SECTION 9: SUMMARY AND NEXT STEPS
% ============================================================================
\section{Summary and Next Steps}

\subsection{Completed Deliverables (Week 3)}

\begin{itemize}
    \item ✓ Project Proposal: Recipe Recommendation and Ingredient Network System.
    \item ✓ Source Inventory: Kaggle recipe dataset with license and access details.
    \item ✓ Schema Draft: Five normalized tables (recipes, ingredients, recipe\_ingredients, ratings, and planned interactions).
    \item ✓ Processed Dataset V1: Ingestion pipeline design and expected output structure.
    \item ✓ Data Dictionary: Complete metadata for all four main tables.
    \item ✓ Scale Analysis: Volume estimates, sparsity metrics, and growth projections.
    \item ✓ Ethics and Access Note: Justification for data use and risk mitigation.
    \item ✓ Reproducibility Specification: Command structure, repository layout, and dependencies.
\end{itemize}

\subsection{Approval Checklist}

The team believes this proposal satisfies the instructor's approval criteria:

\begin{enumerate}
    \item[\checkmark] \textbf{Non-trivial dataset:} 200k+ recipes, 10k+ ingredients, 1M+ interactions with rich metadata and multi-year historical data.
    \item[\checkmark] \textbf{Valid data source:} Public Kaggle Food.com dataset, CC0 license, no unauthorized scraping.
    \item[\checkmark] \textbf{Second-half feasibility:} Schema supports features, interactions, clustering, recommendation, and graph analysis.
    \item[\checkmark] \textbf{Credible build plan:} Reproducible ingestion pipeline, clear joins, and documented cleaning logic.
\end{enumerate}

\subsection{Milestones and Timeline}

\begin{table}[H]
    \centering
    \begin{tabular}{|p{1.5cm}|p{3cm}|p{7.5cm}|}
        \hline
        \textbf{Week} & \textbf{Milestone} & \textbf{Key Deliverables} \\
        \hline
        3 & Dataset Charter V1 & \textit{This report:} schema, ingestion, data dictionary \\
        \hline
        5 & Features \& Dimensionality & PCA/SVD on recipe-ingredient matrix, feature representations \\
        \hline
        7 & Clustering & K-means and DBSCAN analysis, cluster profiles \\
        \hline
        10 & Recommendation Engine & Baseline + collaborative filtering, evaluation metrics \\
        \hline
        12 & Graph Analytics & Ingredient network, centrality, link prediction \\
        \hline
        14 & Final Defense & Integrated report, reproducible repo, live demo \\
        \hline
    \end{tabular}
    \caption{Semester Timeline}
    \label{tab:timeline}
\end{table}

\subsection{Contingency Plans}

Should the primary Kaggle dataset become unavailable:

\begin{enumerate}
    \item Pivot to Food.com or AllRecipes research releases.
    \item Use Recipe1M benchmark dataset (multilingual, larger scale).
    \item Build a smaller custom dataset from open APIs if time permits.
\end{enumerate}

\newpage

% ============================================================================
% REFERENCES
% ============================================================================
\section*{References}

\begin{enumerate}
    \item Kaggle. ``Food Recipes and Reviews Dataset.'' \url{https://www.kaggle.com/datasets/irkaal/food-recipes-and-reviews}. Accessed 2024.
    
    \item Creative Commons. ``CC0 1.0 Universal (Public Domain Dedication).'' \url{https://creativecommons.org/publicdomain/zero/1.0/}. 2024.
    
    \item Jolliffe, I. T., \& Cadima, J. (2016). Principal component analysis: a review and recent developments. \textit{Philosophical Transactions of the Royal Society A}, 374(2065), 20150202.
    
    \item McInnes, L., Healy, J., \& Melville, J. (2018). UMAP: Uniform Manifold Approximation and Projection for Dimension Reduction. \textit{arXiv preprint arXiv:1802.03426}.
    
    \item Koren, Y., Bell, R., \& Volinsky, C. (2009). Matrix Factorization Techniques for Recommender Systems. \textit{Computer}, 42(8), 30--37.
    
    \item Newman, M. E. J. (2010). Networks: An Introduction. Oxford University Press.
    
    \item Bird, S., Klein, E., \& Loper, E. (2009). Natural Language Processing with Python. O'Reilly Media.
\end{enumerate}

\newpage

% ============================================================================
% APPENDIX
% ============================================================================
\appendix

\section{Sample Ingestion Script (Python Pseudocode)}

\begin{lstlisting}[language=Python]
import pandas as pd
import hashlib
from pathlib import Path

def load_raw_data(raw_dir):
    """Load all raw CSV files."""
    recipes = pd.read_csv(f"{raw_dir}/recipes.csv")
    ingredients = pd.read_csv(f"{raw_dir}/ingredients.csv")
    recipe_ingredients = pd.read_csv(f"{raw_dir}/recipe_ingredients.csv")
    ratings = pd.read_csv(f"{raw_dir}/ratings.csv")
    return recipes, ingredients, recipe_ingredients, ratings

def clean_recipes(recipes):
    """Clean recipe table."""
    # Remove duplicates
    recipes = recipes.drop_duplicates(subset=['title', 'instructions'])
    
    # Standardize columns
    recipes['cuisine'] = recipes['cuisine'].fillna('unspecified').str.lower()
    recipes['difficulty'] = recipes['difficulty'].fillna('medium')
    recipes['cook_time'] = recipes['cook_time'].fillna(
        recipes.groupby('cuisine')['cook_time'].transform('median')
    )
    
    return recipes

def clean_ingredients(ingredients):
    """Standardize ingredient names."""
    ingredients['name'] = ingredients['name'].str.lower().str.strip()
    ingredients = ingredients.drop_duplicates(subset=['name'])
    return ingredients

def validate_foreign_keys(recipe_ingredients, recipes, ingredients):
    """Check referential integrity."""
    assert recipe_ingredients['recipe_id'].isin(recipes['recipe_id']).all()
    assert recipe_ingredients['ingredient_id'].isin(ingredients['ingredient_id']).all()
    print("✓ Foreign key validation passed")

def main():
    raw_dir = "data/raw"
    processed_dir = "data/processed"
    
    # Load and clean
    recipes, ingredients, recipe_ingredients, ratings = load_raw_data(raw_dir)
    recipes = clean_recipes(recipes)
    ingredients = clean_ingredients(ingredients)
    
    # Validate
    validate_foreign_keys(recipe_ingredients, recipes, ingredients)
    
    # Save
    Path(processed_dir).mkdir(parents=True, exist_ok=True)
    recipes.to_csv(f"{processed_dir}/recipes_clean.csv", index=False)
    ingredients.to_csv(f"{processed_dir}/ingredients_clean.csv", index=False)
    recipe_ingredients.to_csv(f"{processed_dir}/recipe_ingredients.csv", index=False)
    ratings.to_csv(f"{processed_dir}/ratings_clean.csv", index=False)
    
    print(f"✓ Saved {len(recipes)} recipes, {len(ingredients)} ingredients")

if __name__ == "__main__":
    main()
\end{lstlisting}

\section{Example Data Quality Report}

\begin{lstlisting}[language=json, caption=Sample Quality Report Output]
{
  "recipes": {
    "total_rows": 15000,
    "duplicates_removed": 42,
    "missing_cuisine": 2250,
    "missing_cook_time": 1500,
    "valid_rows": 14958
  },
  "ingredients": {
    "total_rows": 5000,
    "duplicates_removed": 15,
    "missing_category": 600,
    "valid_rows": 4985
  },
  "recipe_ingredients": {
    "total_rows": 150000,
    "foreign_key_violations": 0,
    "valid_rows": 150000
  },
  "ratings": {
    "total_rows": 500000,
    "invalid_scores": 12,
    "missing_review_text": 300000,
    "valid_rows": 499988
  },
  "matrices": {
    "recipe_ingredient_density": 0.002,
    "user_recipe_density": 0.044
  }
}
\end{lstlisting}

\end{document}
